\documentclass[main.tex]{subfiles}


\begin{document}

%from pagenumber to up
% \phantomsection 
% \hypertarget{toc}{}

 

 

 
   

\section{Ход лучей в тонкой линзе}


% Во многих задачах на построения хода лучей света через линзу саму линзу можно считать тонкой (толщина линзы много меньше характериных расстояний в задаче). (Далее предполагается, что речь идет о тонких линзах)

Пусть имеются два параллельных пучка света: один падает на тонкую собирающую линзу, другой~---  на тонкую рассеивающую (рис.\,\ref{fig:p184}).

    \begin{figure}[H]
    \centering

    \begin{tikzpicture}[use optics,font=\small,            line cap=round,                             >={Stealth[inset=0pt,round,length=3mm, angle'=20]},vec/.style={>={Stealth[inset=0pt,round,length=3.5mm, angle'=20]}}]


\node[lens,thick,line cap=butt,object height=3cm,focal length=2cm] (L1) at (0,0) {};
\draw (0,0) node[below left,inner sep=.2222em] {$O$} --(3,0) (0,0)--(-3,0) node[above] {ГОО};
\node[inner sep=0mm,label={[inner xsep=.1111em]below left:$F$}] (FL) at (L1.west focus) {\tikz\draw(0,0) -- (0,2mm);};
\node[inner sep=0mm,label={below:$F$}] (FR) at (L1.east focus) {\tikz\draw(0,0) -- (0,2mm);};


%light
\foreach \i in {0,1,...,4}
\draw[Green,decoration={markings,mark=at position {1/2+1.5/20} with {\arrow{>}}},postaction={decorate},yshift={-\i*.5 cm}] (-2,1) --++ (2,0) coordinate (li\i);

\foreach \i in {0,1,...,4}
\draw[Green,decoration={markings,mark=at position {1/2} with {\arrow{>}}},postaction={decorate},shorten >=-.5cm] (li\i) -- (FR.center);


%%%%%%%
%%%%%%%
%%%%%%%


\begin{scope}[xshift=8cm]

\node[lens,thick,lens type=diverging,line cap=butt,object height=3cm,focal length=2cm] (L1) at (0,0) {};
\draw (0,0) node[below left,inner sep=.2222em] {$O$} --(3,0) (0,0)--(-3,0) node[above] {ГОО};
\node[inner sep=0mm,label={[inner xsep=.1111em]below left:$F$}] (FL) at (L1.west focus) {\tikz\draw(0,0) -- (0,2mm);};
\node[inner sep=0mm,label={below:$F$}] (FR) at (L1.east focus) {\tikz\draw(0,0) -- (0,2mm);};


%light
\foreach \i in {0,1,...,4}
\draw[Green,decoration={markings,mark=at position {1/2+1.5/20} with {\arrow{>}}},postaction={decorate},yshift={-\i*.5 cm}] (-2,1) --++ (2,0) coordinate (li\i);

\foreach \i in {0,1,3,4}{
\draw[densely dashed,gray] (FL.center) -- (li\i);
}

\foreach \i in {0,1,...,4}{
\path[semitransparent,red] 
(FL.center) --++ ({atan((1-\i*.5)/2)}:3.6) coordinate (en\i);
}

\foreach \i in {0,1,...,4}{
\draw[Green,decoration={markings,mark=at position {1/2} with {\arrow{>}}},postaction={decorate}] 
(li\i) -- (en\i);
}

\end{scope}




    \end{tikzpicture}
\caption{Фокусировка и рассеяние параллельного пучка света}
\label{fig:p184}
\end{figure}

В случае собирающей линзы (рис.\,\ref{fig:p184}, слева) пучок света, параллельный \emph{главной оптической оси}~ГОО (оси симметрии линзы), после прохождения линзы собирается в ее правом \emph{фокусе}~$F$ (любая линза имеет два фокуса~--- слева и справа на~ГОО; обычно фокусы расположены на одинаковых расстояниях от линзы). Рассеивающая же линза (рис.\,\ref{fig:p184}, справа) преобразует такой пучок света в расходящийся пучок, как бы (мнимо) выходящий из левого фокуса~$F$ линзы. Точки~$O$ на рис.\,\ref{fig:p184} называются \emph{оптическими центрами} линз. Расстояния~$OF$ называются \emph{фокусными расстояниями} (обозначаются~$F$).

Падающие на линзы лучи на рис.\,\ref{fig:p184} удобно называть \emph{<<нормальными>>} (они как бы падают на линзу под \emph{прямым} углом).

\textbf{Оптическая сила} ($D$ [дптр])~--- это характеристика преломляющей способности линзы:
\begin{equation}
    D=\frac{1}{F}.
\end{equation}
% где $F$~--- фокусное расстояние линзы.

Чем больше оптическая сила линзы, тем ближе к линзе фокусируется (возможно, мнимо) пучок света, падающий на нее <<нормально>> (см.~рис.\,\ref{fig:p184}).

\textbf{Правила хода лучей в линзе} можно сформулировать так\footnote{В случае рассеивающей линзы выражение <<идет через>> понимать как <<мнимо идет через>>.}. 
\begin{enumerate}
    \item <<Нормальный>> луч после преломления идет через фокус (рис.\,\ref{fig:p185}, слева).
    \item Луч, идущий в оптический центр, не преломляется (рис.\,\ref{fig:p185}, посередине).
    \item Если луч падает на линзу \emph{наклонно} (<<ненормально>>), то для построения его дальнейшего хода через центр линзы проводят \emph{побочную оптическую ось}~ПОО, параллельную этому лучу, и находят \emph{побочный фокус~$P$}, находящийся в точке пересечения~ПОО с \emph{фокальной плоскостью}~ФП, проходящей через фокус~$F$ перпендикулярно~ГОО. Преломленный луч идет через этот побочный фокус (рис.\,\ref{fig:p185}, справа). 
\end{enumerate}
\nointerlineskip

    \begin{figure}[H]
    \centering

    \begin{tikzpicture}[use optics,font=\small,            line cap=round,                             >={Stealth[inset=0pt,round,length=3mm, angle'=20]},vec/.style={>={Stealth[inset=0pt,round,length=3.5mm, angle'=20]}}]


\node[lens,thick,line cap=butt,object height=2cm,focal length=1cm] (L1) at (0,0) {};
\draw (0,0) node[below left,inner sep=.2222em] {} -- (2,0) (0,0) -- (-2,0);
\node[inner sep=0mm,label={below:$F$}] (FL) at (L1.west focus) {\tikz\draw(0,0) -- (0,2mm);};
\node[inner sep=0mm,label={below:$F$}] (FR) at (L1.east focus) {\tikz\draw(0,0) -- (0,2mm);};

\draw[Green,decoration={markings,mark=at position {1/2+1.5/20} with {\arrow{>}}},postaction={decorate}] (-2,.5) --++ (2,0) coordinate (li);
\draw[Green,decoration={markings,mark=at position {1/2+1.5/20} with {\arrow{>}}},postaction={decorate},shorten >=-.5cm] (li) -- (FR.center);


%%%


\begin{scope}[xshift=5cm]

\node[lens,thick,line cap=butt,object height=2cm,focal length=1cm] (L1) at (0,0) {};
\draw (0,0) node[below left,inner sep=.2222em] {} -- (2,0) (0,0) -- (-2,0);
\node[inner sep=0mm,label={below:$F$}] (FL) at (L1.west focus) {\tikz\draw(0,0) -- (0,2mm);};
\node[inner sep=0mm,label={above:$F$}] (FR) at (L1.east focus) {\tikz\draw(0,0) -- (0,2mm);};

\draw[Green,decoration={markings,mark=at position {1/2+1.5/20} with {\arrow{>}}},postaction={decorate}] (-2,.5) -- (0,0);
\draw[Green,decoration={markings,mark=at position {1/2+1.5/20} with {\arrow{>}}},postaction={decorate}] (0,0) -- (2,-.5);


\end{scope}


%%%


\begin{scope}[xshift=10cm]

\node[lens,thick,line cap=butt,object height=2cm,focal length=1cm] (L1) at (0,0) {};
\draw (0,0) node[below right,inner sep=.1111em] {} -- (2,0) (0,0) -- (-2,0);
\node[inner sep=0mm,label={below:}] (FL) at (L1.west focus) {\tikz\draw(0,0) -- (0,2mm);};
\node[inner sep=0mm,label={[inner xsep=.1111em]below right:$F$}] (FR) at (L1.east focus) {\tikz\draw(0,0) -- (0,2mm);};


\draw[Green,decoration={markings,mark=at position {1/2+1.5/20} with {\arrow{>}}},postaction={decorate}] (-1.5,0) -- (0,.5);
\coordinate (O) at (0,0);
\coordinate (b) at (-1.5,-.5);
\draw[densely dashed, red,name path=red] (b) node[below,black] {ПОО} -- (O) coordinate[pos=2](en) -- (en);
\draw[densely dashed,name path=gr] (FR.center) ++ (0,1) --++ (0,-2) node [right] {ФП};
\draw[Green,decoration={markings,mark=at position {1/2+1.5/20} with {\arrow{>}}},postaction={decorate},name intersections={of=red and gr, by=x},shorten >=-.5cm] (0,.5) -- (x);
\path[name intersections={of=red and gr, by=x}] (x) node[] [circle,fill=,inner sep=0pt,minimum size=1mm,label={[inner xsep=.1111em]above right:$P$}] {};
;


\end{scope}



    \end{tikzpicture}
\caption{К правилам хода лучей в линзе}
\label{fig:p185}
\end{figure}










 
\end{document}
